\documentclass[aps,preprint]{revtex4}%
\usepackage{amsfonts}%
\usepackage{amsmath}%
\setcounter{MaxMatrixCols}{30}%
\usepackage{amssymb}%
\usepackage{graphicx}
%TCIDATA{OutputFilter=latex2.dll}
%TCIDATA{Version=4.10.0.2347}
%TCIDATA{CSTFile=revtex4.cst}
%TCIDATA{Created=Friday, December 20, 2002 11:44:01}
%TCIDATA{LastRevised=Friday, December 20, 2002 17:28:34}
%TCIDATA{<META NAME="GraphicsSave" CONTENT="32">}
%TCIDATA{<META NAME="DocumentShell" CONTENT="Articles\SW\REVTeX 4">}
\newtheorem{theorem}{Theorem}
\newtheorem{acknowledgement}[theorem]{Acknowledgement}
\newtheorem{algorithm}[theorem]{Algorithm}
\newtheorem{axiom}[theorem]{Axiom}
\newtheorem{claim}[theorem]{Claim}
\newtheorem{conclusion}[theorem]{Conclusion}
\newtheorem{condition}[theorem]{Condition}
\newtheorem{conjecture}[theorem]{Conjecture}
\newtheorem{corollary}[theorem]{Corollary}
\newtheorem{criterion}[theorem]{Criterion}
\newtheorem{definition}[theorem]{Definition}
\newtheorem{example}[theorem]{Example}
\newtheorem{exercise}[theorem]{Exercise}
\newtheorem{lemma}[theorem]{Lemma}
\newtheorem{notation}[theorem]{Notation}
\newtheorem{problem}[theorem]{Problem}
\newtheorem{proposition}[theorem]{Proposition}
\newtheorem{remark}[theorem]{Remark}
\newtheorem{solution}[theorem]{Solution}
\newtheorem{summary}[theorem]{Summary}
\newenvironment{proof}[1][Proof]{\noindent\textbf{#1.} }{\ \rule{0.5em}{0.5em}}
\begin{document}
\preprint{HEP/123-qed}
\title[Short title for running header]{Long title that appears on the title page}
\author{First Author}
\affiliation{My Institution}
\author{Second Author}
\affiliation{My Institution}
\author{Third Author}
\affiliation{Other Institution}
\keywords{one two three}
\pacs{PACS number}

\begin{abstract}
Shell document for REV\TeX{} 4.

\end{abstract}
\volumeyear{year}
\volumenumber{number}
\issuenumber{number}
\eid{identifier}
\date[Date text]{date}
\received[Received text]{date}

\revised[Revised text]{date}

\accepted[Accepted text]{date}

\published[Published text]{date}

\startpage{101}
\endpage{102}
\maketitle
\tableofcontents


\section{Transition Amplitudes and Expectation Values wrt IPSs}

We will need to calculate expectation values and transition amplitudes between
IPS\ composed of non orthogonal 1-particle states, we will follow the
presentation found in "On the Evaluation of Nonorthogonal Matrix
Elements\textit{" }by Jacob Verbeek and J.H. van Lenthe THEOCHEM 1990.
Consider two IPSs $\left\vert \Psi\right\rangle $ and $\left\vert
\Phi\right\rangle $ composed of nonorthogonal general spin orbitals $\left\{
\psi_{j}|1\leq j\leq N\right\}  $ and $\left\{  \varphi_{j}|1\leq j\leq
N\right\}  $ and the transition amplitude
\begin{equation}
\left\langle \Psi|H\Phi\right\rangle
\end{equation}
where the Hamiltonian is a sum of one and two particle operators
\begin{equation}
H=%
%TCIMACRO{\dsum \limits_{1\leq j\leq N}}%
%BeginExpansion
{\displaystyle\sum\limits_{1\leq j\leq N}}
%EndExpansion
h\left(  j\right)  +\frac{1}{2}%
%TCIMACRO{\dsum \limits_{1\leq j\neq k\leq N}}%
%BeginExpansion
{\displaystyle\sum\limits_{1\leq j\neq k\leq N}}
%EndExpansion
g\left(  j,k\right)
\end{equation}


The transition amplitude of the $1$-particle part is given by
\begin{equation}
\left\langle \Psi\left\vert
%TCIMACRO{\dsum \limits_{1\leq j\leq N}}%
%BeginExpansion
{\displaystyle\sum\limits_{1\leq j\leq N}}
%EndExpansion
h\left(  j\right)  \right.  \Phi\right\rangle =Tr\left\{  \mathbf{h}%
\operatorname{adj}\left(  \mathbf{S}\right)  \right\}
\end{equation}
where the matrix $\mathbf{h}$ has matrix elements $\left\{  h_{jk}|1\leq
j,k\leq N\right\}  $ given by
\begin{equation}
h_{jk}=\left\langle \psi_{j}|h\left(  1\right)  \varphi_{k}\right\rangle =%
%TCIMACRO{\dint }%
%BeginExpansion
{\displaystyle\int}
%EndExpansion
\psi_{j}\left(  \mathsf{z}\right)  ^{\ast}h\left(  \mathsf{z,z}^{\prime
}\right)  \varphi_{k}\left(  \mathsf{z}^{\prime}\right)  d\mathsf{z}%
d\mathsf{z}^{\prime}%
\end{equation}
and the classical $N\times N$ adjugate matrix $\operatorname{adj}\left(
\mathbf{S}\right)  $ of the overlap matrix $\mathbf{S\equiv}\left\{
\left\langle \psi_{j}|\varphi_{k}\right\rangle |1\leq j,k\leq N\right\}  $ is
defined by
\begin{equation}
\left(  \operatorname{adj}\left(  \mathbf{S}\right)  \right)  _{jk}=\left(
-1\right)  ^{j+k}\det\left(  \mathbf{S}\left[  k|j\right]  \right)
\end{equation}
where $\mathbf{S}\left[  k|j\right]  $ is the $\left(  N-1\right)  \times$
$\left(  N-1\right)  $ matrix obtained from $\mathbf{S}$ by deleting the row
$k$ and the column $j$.

The transition amplitude of the $2$-particle part is given by
\begin{equation}
\left\langle \Psi\left\vert \frac{1}{2}%
%TCIMACRO{\dsum \limits_{1\leq j\neq k\leq N}}%
%BeginExpansion
{\displaystyle\sum\limits_{1\leq j\neq k\leq N}}
%EndExpansion
g\left(  j,k\right)  \right.  \Phi\right\rangle =Tr\left\{  \mathbf{G}%
\operatorname{adj}^{\left(  2\right)  }\left(  \mathbf{S}\right)  \right\}
\end{equation}
where the matrix $\mathbf{G}$ has matrix elements $\left\{  g_{jklm}|1\leq
j<k\leq N;1\leq l<m\leq N\right\}  $ given by
\begin{align}
g_{jklm}  & =\left\langle \psi_{j}\psi_{k}|g\left(  j,k\right)  \varphi
_{l}\varphi_{m}\right\rangle \nonumber\\
& =%
%TCIMACRO{\dint }%
%BeginExpansion
{\displaystyle\int}
%EndExpansion
\left\{  \psi_{j}\left(  \mathsf{z}_{1}\right)  ^{\ast}\psi_{k}\left(
\mathsf{z}_{2}\right)  ^{\ast}g\left(  \mathsf{z}_{1}\mathsf{,\mathsf{z}%
_{2},\mathsf{z}_{1}^{\prime},z}_{2}^{\prime}\right)  \varphi_{l}\left(
\mathsf{z}_{1}^{\prime}\right)  \varphi_{m}\left(  \mathsf{z}_{2}^{\prime
}\right)  \right.  \nonumber\\
& \left.  -\psi_{j}\left(  \mathsf{z}_{1}\right)  ^{\ast}\psi_{k}\left(
\mathsf{z}_{2}\right)  ^{\ast}g\left(  \mathsf{z}_{1}\mathsf{,\mathsf{z}%
_{2},\mathsf{z}_{1}^{\prime},z}_{2}^{\prime}\right)  \varphi_{m}\left(
\mathsf{z}_{1}^{\prime}\right)  \varphi_{l}\left(  \mathsf{z}_{2}^{\prime
}\right)  \right\}  d\mathsf{z}^{2}d\mathsf{z}^{\prime2}%
\end{align}
and the second order $\binom{N}{2}\times$ $\binom{N}{2}$ adjugate matrix
$\operatorname{adj}^{\left(  2\right)  }\left(  \mathbf{S}\right)  $ is
defined by
\begin{equation}
\left(  \operatorname{adj}^{\left(  2\right)  }\left(  \mathbf{S}\right)
\right)  _{jklm}=\left(  -1\right)  ^{j+k}\det\left(  \mathbf{S}\left[
lm|jk\right]  \right)
\end{equation}
where

$\mathbf{S}\left[  lm|jk\right]  $ is the $\left(  N-2\right)  \times\left(
N-2\right)  $ matrix obtained from $\mathbf{S}$ by deleting the rows $l,m$ and
the column $j,k$.




\end{document}